The distance map \(z\mapsto d(z,x_j)\) is one-Lipschitz. Thus \(\psi_{j,h}(z)=(1-d(z,x_j)/h)_+\) is \(h^{-1}\)-Lipschitz, and each perturbed side regression in \(\mathrm{def:fano\text{-}tent\text{-}family}\) is \(L/4\)-Lipschitz. Its absolute trace is at most \(Lh/4\le L\) because \(0<h\le h_0/6<1\). The perturbation changes neither the score law nor any local-mass or packing property. The sharp-assignment and consistency properties follow directly from the construction.

For each Fano law \(P\) and side \(t\in\{0,1\}\), set
\[
 N_{t,P}=\mathcal X_{t,P}\setminus\mathcal A_t.
\]
The set \(\mathcal X_{t,P}\) is closed by its definition as a support and
\(\mathcal A_t\) is Borel, so \(N_{t,P}\) is Borel.  Moreover, for the
one-sided restricted score measure,
\[
 \mu_{t,P}(N_{t,P})
 =P_X\bigl((\mathcal X_{t,P}\setminus\mathcal A_t)
             \cap\mathcal A_t\bigr)=0. \tag{a}
\]
In the frozen circular construction,
\[
 \mathcal A_1=\{z:\rho(z)\le0\},
 \qquad
 \mathcal A_0=\{z:\rho(z)>0\}
 =\mathbb R^2\setminus\mathcal A_1.
\]
Hence \(z\in\mathcal X_{t,P}\setminus N_{t,P}\) implies
\(z\in\mathcal A_t\), and the sharp assignment rule gives \(T(z)=t\).
Using the displayed tent formulas as Borel extensions off their
one-sided supports, a single regular conditional kernel for the
observed outcome is
\[
 \mathsf R_P(z,dy)
 =\mathcal N(g_{T(z),P}(z),\sigma^2)(dy).
\]
For every \(z\in\mathcal X_{t,P}\setminus N_{t,P}\), this is
\(\mathcal N(g_{t,P}(z),\sigma^2)\).  Therefore
\[
 \int_{\mathbb R}|y-g_{t,P}(z)|\,\mathsf R_P(z,dy)<\infty,
 \qquad
 \int_{\mathbb R}\{y-g_{t,P}(z)\}\,\mathsf R_P(z,dy)=0,
\]
and, for every \(\lambda\in\mathbb R\),
\[
 \int_{\mathbb R}e^{\lambda\{y-g_{t,P}(z)\}}
 \,\mathsf R_P(z,dy)
 =e^{\sigma^2\lambda^2/2}.
\]
Thus the integrability, centering, and moment-generating-function
requirements of \(\mathrm{ass:gaussian\text{-}errors}\) hold with the
side-null sets in (a).  Boundary points assigned to the opposite side
are absorbed into \(N_{t,P}\).  The nullity calculation (a) itself does
not require overlap, while the implication \(T(z)=t\) here uses the
explicit complementary circular partition.  Consequently
\(\mathrm{lem:circular\text{-}witness\text{-}membership}\) and the class
definitions put every pervasive perturbation in
\(\mathcal P_{\kappa}^{\mathrm{perv}}\) and both isolated perturbations in
\(\mathcal P_{\kappa}^{\mathrm{iso}}\).

Maximality of the \(3h\)-separated set implies that arcs of radius \(3h\) cover \(S^1\), while arcs of radius \(3h/2\) about its points are disjoint. Comparing their total arclength with \(2\pi\) gives
\[
 c h^{-1}\le |\Xi_h^{\mathrm{perv}}|\le C h^{-1}. \tag{7}
\]

Put \(a=Lh/4\). Conditional on \(X=z\), the relevant Gaussian coordinates under alternatives \(j\) and \(k\) have common variance \(\sigma^2\), and their mean difference has absolute value \(a|\psi_{j,h}(z)-\psi_{k,h}(z)|\). Direct integration of the Gaussian density gives
\[
 \operatorname{KL}\{\mathcal N(\mu,\sigma^2),\mathcal N(\nu,\sigma^2)\}
 =\frac{(\mu-\nu)^2}{2\sigma^2}. \tag{8}
\]
The integrand vanishes outside \(B(x_j,h)\cup B(x_k,h)\). The coordinate calculations in \(\mathrm{lem:circular\text{-}witness\text{-}membership}\) give
\[
 \int_{B(x_j,h)}f_{\kappa}^{\mathrm{perv}}(z)\,dz\le Ch^{\kappa},
 \qquad
 \int_{B(x_{\star},h)}f_{\kappa}^{\mathrm{iso}}(z)\,dz\le Ch^{\kappa}.
\]
Since \(0\le\psi_{j,h}\le1\), conditioning on \(X\) in (8), and summing the two potential-outcome coordinates if the superpopulation experiment is used, gives in either family
\[
 \operatorname{KL}(\mathbb P_{\kappa,j,h},\mathbb P_{\kappa,k,h})
 \le \frac{Ca^2h^{\kappa}}{\sigma^2}
 \le \frac{CL^2h^{\kappa+2}}{\sigma^2}. \tag{9}
\]
The observed-data divergence is no larger, by marginalization. For the displayed product measures, likelihood ratios factor by the definition of a product law, so relative entropy is additive. Thus
\[
 \operatorname{KL}(\mathbb P_{\kappa,j,h}^{\otimes n},
 \mathbb P_{\kappa,k,h}^{\otimes n})
 \le \frac{CnL^2h^{\kappa+2}}{\sigma^2}. \tag{10}
\]
At an active center, \(\psi_{j,h}(x_j)=1\). The trace formulas in \(\mathrm{def:fano\text{-}tent\text{-}family}\) yield
\[
 \theta_{1,P}(x_j)=\frac{Lh}{4},\qquad
 \theta_{0,P}(x_j)=-\frac{Lh}{4},\qquad
 \tau_P(x_j)=\frac{Lh}{2}.
\]
Together with (7) and (10), these are all the asserted properties.